
\documentclass{article}
\usepackage{colortbl}
\usepackage{makecell}
\usepackage{multirow}
\usepackage{supertabular}

\begin{document}

\newcounter{utterance}

\centering \large Interaction Transcript for game `eqbench', experiment `eq\_bench', episode 111 with reflexion{-}Qwen3.5{-}27B.
\vspace{24pt}

{ \footnotesize  \setcounter{utterance}{1}
\setlength{\tabcolsep}{0pt}
\begin{supertabular}{c@{$\;$}|p{.15\linewidth}@{}p{.15\linewidth}p{.15\linewidth}p{.15\linewidth}p{.15\linewidth}p{.15\linewidth}}
    \# & \multicolumn{2}{c}{Player} && \multicolumn{2}{c}{Game Master} \\
    \hline

    \theutterance \stepcounter{utterance}  
    & & & \multicolumn{4}{p{0.6\linewidth}}{
        \cellcolor[rgb]{0.9,0.9,0.9}{
            \makecell[{{p{\linewidth}}}]{
                \texttt{\tiny{[P1$\langle$GM]}}
                \texttt{You have attempted this same task before and did not succeed. These are the plans you wrote after those attempts.} \\
\texttt{Trial \#1: New plan: My previous attempt failed because I assigned a high score (8) to "Humiliated," which contradicts the character's behavior. Luca ends the dialogue by aggressively rejecting Charlie's authority ("You can keep your stripes. I don't need them"), which is a classic defense mechanism to protect his ego, not a sign of feeling humiliated. A humiliated person would likely withdraw or show shame, whereas Luca remains combative and self{-}justifying. Next time, I will analyze the final lines specifically for signs of ego protection versus shame. I will assign a score of 0 to "Humiliated" and "Complacent," a low score to "Relieved" (as the conflict is unresolved), and a high score (e.g., 9 or 10) to "Defensive," as the entire dialogue is a reaction to criticism where Luca doubles down on his status and attacks the critic's character.} \\
\\ 
\texttt{Trial \#2: New plan: My previous attempt failed because I assigned a non{-}zero score (1) to "Relieved," which contradicts the character's state. Luca ends the dialogue in a state of high tension and unresolved conflict, having just been attacked personally and professionally; there is no indication of relief or release of tension. To correct this, I will strictly evaluate the final lines for evidence of tension release. Since Luca is actively rejecting the critic and doubling down on his ego ("I don't need them"), he is clearly not relieved. I will assign a score of 0 to "Relieved," "Humiliated," and "Complacent," and reserve the maximum score (10) exclusively for "Defensive," as this is the only emotion supported by his aggressive, self{-}justifying, and combative closing statements.} \\
\\ 
\\ 
\\ 
\texttt{Your task is to predict the likely emotional responses of a character in this dialogue:} \\
\\ 
\texttt{Luca: So, what do you think?} \\
\texttt{Charlie: It's... interesting.} \\
\texttt{Luca: Interesting? That's it?} \\
\texttt{Charlie: Well, it's not exactly a compliment, is it?} \\
\texttt{Luca: You haven't even tasted it yet. } \\
\texttt{Charlie: I don't need to taste it to know it's a mess.} \\
\texttt{Luca: A mess? You're calling my masterpiece a mess?} \\
\texttt{Charlie: If the shoe fits...} \\
\texttt{Luca: You're just a critic, Charlie. You don't know what it takes to create something.} \\
\texttt{Charlie: And you don't know what it takes to appreciate something. } \\
\texttt{Luca: I've been a chef for over 20 years. I've earned my stars.} \\
\texttt{Charlie: And I've been a critic for 15. I've earned my stripes.} \\
\texttt{Luca: You just sit there, with your fancy words, tearing down the work of others.} \\
\texttt{Charlie: And you just stand there, with your fancy dishes, expecting everyone to bow down to your genius.} \\
\texttt{Luca: You're just bitter because your life is falling apart.} \\
\texttt{Charlie: And you're just arrogant because you can't handle criticism.} \\
\texttt{Luca: I don't need to take this from you.} \\
\texttt{Charlie: And I don't need to eat this from you.} \\
\texttt{Luca: You know what, Charlie? You can keep your stripes. I don't need them.} \\
\texttt{Charlie: Good. Because I wasn't planning on giving them to you anyway.} \\
\texttt{{[}End dialogue{]}} \\
\\ 
\texttt{At the end of this dialogue, Luca would feel...} \\
\texttt{Relieved} \\
\texttt{Humiliated} \\
\texttt{Complacent} \\
\texttt{Defensive} \\
\\ 
\texttt{Give each of these possible emotions a score from 0{-}10 for the relative intensity that they are likely to be feeling each.} \\
\\ 
\texttt{You must output in the following format, including headings (of course, you should give your own scores), with no additional commentary:} \\
\\ 
\texttt{Relieved: <score>} \\
\texttt{Humiliated: <score>} \\
\texttt{Complacent: <score>} \\
\texttt{Defensive: <score>} \\
\\ 
\\ 
\texttt{{[}End of answer{]}} \\
\\ 
\texttt{Remember: zero is a valid score, meaning they are likely not feeling that emotion. You must score at least one emotion > 0.} \\
\\ 
\texttt{Your answer:} \\
            }
        }
    }
     \\ \\

    \theutterance \stepcounter{utterance}  
    & \multicolumn{4}{p{0.6\linewidth}}{
        \cellcolor[rgb]{0.561,0.737,0.561}{
            \makecell[{{p{\linewidth}}}]{
                \texttt{\tiny{[P1$\rangle$GM]}}
                \texttt{Relieved: 0} \\
\texttt{Humiliated: 0} \\
\texttt{Complacent: 0} \\
\texttt{Defensive: 10} \\
            }
        }
    }
    & & \\ \\

    \theutterance \stepcounter{utterance}  
    & & & \multicolumn{2}{p{0.3\linewidth}}{
        \cellcolor[rgb]{0.9,0.9,0.9}{
            \makecell[{{p{\linewidth}}}]{
                \texttt{\tiny{[GM$|$GM]}}
                \texttt{Relieved: 0} \\
\texttt{Humiliated: 0} \\
\texttt{Complacent: 0} \\
\texttt{Defensive: 10} \\
            }
        }
    }
    & & \\ \\

    \theutterance \stepcounter{utterance}  
    & & & \multicolumn{2}{p{0.3\linewidth}}{
        \cellcolor[rgb]{0.9,0.9,0.9}{
            \makecell[{{p{\linewidth}}}]{
                \texttt{\tiny{[GM$|$GM]}}
                \texttt{Relieved: 0} \\
\texttt{Humiliated: 6} \\
\texttt{Complacent: 0} \\
\texttt{Defensive: 8} \\
            }
        }
    }
    & & \\ \\

    \theutterance \stepcounter{utterance}  
    & & & \multicolumn{2}{p{0.3\linewidth}}{
        \cellcolor[rgb]{0.9,0.9,0.9}{
            \makecell[{{p{\linewidth}}}]{
                \texttt{\tiny{[GM$|$GM]}}
                \texttt{game\_result = LOSE} \\
            }
        }
    }
    & & \\ \\

    \theutterance \stepcounter{utterance}  
    & & \multicolumn{4}{p{0.6\linewidth}}{
        \cellcolor[rgb]{0.95,0.95,0.95}{
            \makecell[{{p{\linewidth}}}]{
                \texttt{\tiny{[P1$\leftrightarrow$P]}}
                \texttt{Post{-}episode Reflexion (not sent to GM):} \\
\texttt{New plan: My previous attempt failed because I assigned a score of 0 to "Humiliated," overlooking the underlying emotional reality of the character. While Luca's *behavior* is defensive and aggressive, the *cause* of this behavior is the humiliation of having his "masterpiece" and 20 years of experience dismissed as a "mess" by a critic. His final line ("I don't need them") is a desperate attempt to mask the sting of rejection, which is a hallmark of humiliation. A purely defensive character might simply argue facts; a humiliated character attacks the source of the pain to protect their ego. Next time, I will recognize that "Defensive" describes the *action* while "Humiliated" describes the *feeling* driving that action. I will assign a high score (e.g., 9) to "Defensive" for his combative stance, but also a significant score (e.g., 7 or 8) to "Humiliated" to reflect the deep personal insult he is reacting to, while keeping "Relieved" and "Complacent" at 0.} \\
            }
        }
    }
    & \\ \\

\end{supertabular}
}

\end{document}
